\medskip

\textbf{Principaux domaines abordés :} géométrie dans l'espace.

\medskip

Une maison est constituée d'un parallélépipède rectangle ABCDEFGH surmonté d'un prisme EFIHGJ dont une base est le triangle EIF isocèle en I.

Cette maison est représentée ci-dessous.

\begin{center}
\psset{unit=1cm,arrowsize=2pt 3}
\begin{pspicture}(12.2,6)
\pspolygon(0.4,3.4)(0.4,1.5)(1.9,0.4)(1.9,2.3)
\psline(1.9,0.4)(7.9,0.8)(7.9,2.7)(1.9,2.3)
\psline[linestyle=dashed](0.4,1.5)(6.4,1.9)(6.4,3.8)(0.4,3.4)
\psline[linestyle=dashed](7.9,0.8)(6.4,1.9)
\pspolygon(7.9,2.7)(6.4,3.8)(3.4,5.5)(4.9,4.4)
\psline(4.9,4.4)(1.9,2.3)
\psline(3.4,5.5)(0.4,3.4)
\psline[linewidth=1.5pt]{->}(1.9,0.4)(3.9,0.533)
\psline[linewidth=1.5pt]{->}(1.9,0.4)(1.15,0.95)
\psline[linewidth=1.5pt]{->}(1.9,0.4)(1.9,2.3)
\psline(11.6,0.3)(7.9,1.2)
\psline[linestyle=dashed](7.9,1.2)(7.15,1.35)
\uput[d](1.9,0.4){A} \uput[dr](7.9,0.8){B} \uput[r](6.4,1.9){C} \uput[l](0.4,1.5){D}
\uput[u](1.9,2.3){E} \uput[r](7.9,2.7){F} \uput[ur](6.4,3.8){G} \uput[ul](0.4,3.4){H}
\uput[l](4.9,4.4){I} \uput[u](3.4,5.5){J} \uput[r](11.6,0.3){R} \uput[dl](1.7,0.6){$\vect{\jmath}$}
\uput[dr](2.9,0.45){$\vect{\imath}$} \uput[r](1.9,1.35){$\vect{k}$}
\psdots(11.6,0.3)
\end{pspicture}
\end{center}

On a AB $= 3$,\quad AD $= 2$,\quad AE $= 1$.

On définit les vecteurs $\vect{\imath}= \dfrac13\vect{\text{AB}},\:\vect{\jmath}= \dfrac12\vect{\text{AD}}, \:\vect{k} = \vect{\text{AE}}$.

On munit ainsi l'espace du repère orthonormé $\left(\text{A}~;~\vect{\imath},~\vect{\jmath},~\vect{k}\right)$.

\medskip

\begin{enumerate}
\item Donner les coordonnées du point G.
\item Le vecteur $\vect{n}$ de coordonnées $(2~;~0~;~-3)$ est vecteur normal au plan (EHI).

Déterminer une équation cartésienne du plan (EHI).
\item Déterminer les coordonnées du point I.
\item Déterminer une mesure au degré près de l'angle $\widehat{\text{EIF}}$.
\item Afin de raccorder la maison au réseau électrique, on souhaite creuser une tranchée rectiligne depuis un relais électrique situé en contrebas de la maison.

Le relais est représenté par le point R de coordonnées $(6~;~- 3~;~- 1)$.

La tranchée est assimilée à un segment d'une droite $\Delta$ passant par R et dirigée par le vecteur $\vect{u}$ de coordonnées $(-3~;~4~;~1)$. On souhaite vérifier que la tranchée atteindra la maison au niveau de l'arête [BC].
	\begin{enumerate}
		\item Donner une représentation paramétrique de la droite $\Delta$.
		\item On admet qu'une équation du plan (BFG) est $x = 3$.
		
Soit K le point d'intersection de la droite $\Delta$ avec le plan (BFG).

Déterminer les coordonnées du point K.
		\item Le point K appartient-il bien à l'arête [BC] ?
	\end{enumerate}
\end{enumerate}

\bigskip
